% Source: Evans, "Partial Differential Equations" (2nd ed., AMS Graduate Studies in Mathematics vol. 19, 2010),
% Chapter C "Calculus Facts", PDF pages 628-636 of MathBooks/Evans Partial Differential Equations(1).pdf.
% Transcribed page-by-page by vision subagents, then merged and restructured into
% leanblueprint conventions (semantic \label, bracketed titles, \uses dependency edges) below.

\section{Calculus Facts}
\label{sec:calculus-facts}

\subsection{Boundaries}

Let $U \subset \mathbb{R}^n$ be open and bounded, $k \in \{1,2,\ldots\}$.

\begin{definition}[$C^k$ boundary]
\label{def:ck-boundary-regularity}
\dcref{appC:C.1-def-1}
\group{Boundary Regularity}
We say $\partial U$ is $C^k$ if for each point $x^0 \in \partial U$ there exist $r > 0$ and a $C^k$ function $\gamma : \mathbb{R}^{n-1} \to \mathbb{R}$ such that---upon relabeling and reorienting the coordinates axes if necessary---we have
$$U \cap B(x^0, r) = \{x \in B(x^0, r) \mid x_n > \gamma(x_1, \ldots, x_{n-1})\}.$$

Likewise, $\partial U$ is $C^\infty$ if $\partial U$ is $C^k$ for $k = 1, 2, \ldots$, and $\partial U$ is analytic if the mapping $\gamma$ is analytic.
\end{definition}

\begin{figure}[h]\centering
% [FIGURE: A diagram showing the boundary of region U. The figure depicts a curved boundary between a shaded region above (exterior) and an unshaded region below (interior). There is a coordinate system at the bottom labeled $\mathbb{R}^{n-1}$, and annotations showing $\gamma$, a point $x^0$ on the boundary, the outward normal vector $\nu(x^0)$, and the boundary itself labeled $\partial U$.]
\end{figure}

The boundary of $U$

\begin{definition}[Outward normal and normal derivative]
\label{def:outward-normal-and-normal-derivative}
\dcref{appC:C.1-def-2}
\group{Boundary Regularity}
\uses{def:ck-boundary-regularity}
(i) If $\partial U$ is $C^1$, then along $\partial U$ is defined the outward pointing unit normal vector field
$$\nu = (\nu^1, \ldots, \nu^n).$$

The unit normal at any point $x^0 \in \partial U$ is $\nu(x^0) = \nu = (\nu_1, \ldots, \nu_n)$.

(ii) Let $u \in C^1(\overline{U})$. We call
$$\frac{\partial u}{\partial \nu} := \nu \cdot Du$$
the (outward) normal derivative of $u$.
\end{definition}

We will frequently need to change coordinates near a point of $\partial U$ so as to ``flatten out'' the boundary. To be more specific, fix $x^0 \in \partial U$, and choose $r, \gamma$, etc.\ as above. Define then
$$\begin{cases} y_i = x_i =: \Phi^i(x) & (i = 1, \ldots, n-1) \\ y_n = x_n - \gamma(x_1, \ldots, x_{n-1}) =: \Phi^n(x), \end{cases}$$
and write
$$y = \mathbf{\Phi}(x).$$

Similarly, we set
$$\begin{cases} x_i = y_i =: \Psi^i(y) & (i = 1, \ldots, n-1) \\ x_n = y_n + \gamma(y_1, \ldots, y_{n-1}) =: \Psi^n(y), \end{cases}$$
and write
$$x = \mathbf{\Psi}(y).$$

Then $\mathbf{\Phi} = \mathbf{\Psi}^{-1}$, and the mapping $x \mapsto \mathbf{\Phi}(x) = y$ ``straightens out $\partial U$'' near $x^0$. Observe also that $\det D\mathbf{\Phi} = \det D\mathbf{\Psi} = 1$.

\begin{figure}[h]\centering
% [FIGURE: Two disks joined by curved arrows labeled $\Phi$ (top, pointing right) and $\Psi$ (bottom, pointing left). The left disk, labeled "x-coordinates", shows a shaded region above a curved boundary line through a marked point $x^0$. The right disk, labeled "y-coordinates", shows the corresponding shaded region above a straight horizontal boundary line through a marked point $y^0$, illustrating how the map $\Phi$ flattens the curved boundary into a flat one.]
\end{figure}

Straightening out the boundary

\subsection{Gauss-Green Theorem}

In this section we assume $U$ is a bounded, open subset of $\mathbb{R}^n$, and $\partial U$ is $C^1$.

\begin{theorem}[Gauss-Green Theorem]
\label{thm:gauss-green-theorem}
\dcref{appC:thm-1}
\group{Gauss-Green Theorem}
\uses{def:ck-boundary-regularity,def:outward-normal-and-normal-derivative}
Suppose $u \in C^1(\overline{U})$. Then
$$(1) \quad \int_U u_{x_i} \, dx = \int_{\partial U} u \nu^i \, dS \quad (i = 1, \ldots, n).$$
\end{theorem}

\begin{theorem}[Integration-by-parts formula]
\label{thm:integration-by-parts-formula}
\dcref{appC:thm-2}
\group{Gauss-Green Theorem}
\uses{thm:gauss-green-theorem}
Let $u, v \in C^1(\overline{U})$. Then
$$(2) \quad \int_U u_{x_i} v \, dx = -\int_U u v_{x_i} \, dx + \int_{\partial U} u v \nu^i \, dS \quad (i = 1, \ldots, n).$$
\end{theorem}

\begin{proof}
Apply \cref{thm:gauss-green-theorem} to $uv$.
\end{proof}

\begin{theorem}[Green's formulas]
\label{thm:greens-formulas}
\dcref{appC:thm-3}
\group{Gauss-Green Theorem}
\uses{thm:integration-by-parts-formula,def:outward-normal-and-normal-derivative}
Let $u, v \in C^2(\overline{U})$. Then

(i) $$\int_U \Delta u \, dx = \int_{\partial U} \frac{\partial u}{\partial \nu} \, dS,$$

(ii) $$\int_U Dv \cdot Du \, dx = -\int_U u \Delta v \, dx + \int_{\partial U} \frac{\partial v}{\partial \nu} u \, dS,$$

(iii) $$\int_U u \Delta v - v \Delta u \, dx = \int_{\partial U} u \frac{\partial v}{\partial \nu} - v \frac{\partial u}{\partial \nu} \, dS.$$
\end{theorem}

\begin{proof}
Using (2), with $u_{x_i}$ in place of $u$ and $v \equiv 1$, we see
$$\int_U u_{x_i x_i} \, dx = \int_{\partial U} u_{x_i} \nu^i \, dS.$$

Sum $i = 1, \ldots, n$ to establish (i).

To derive (ii), we employ (2) with $v = u_{x_i}$. Write (ii) with $u$ and $v$ interchanged and then subtract, to obtain (iii).
\end{proof}

\subsection{Polar Coordinates, Coarea Formula}

Next we convert $n$-dimensional integrals into integrals over spheres.

\begin{theorem}[Polar coordinates]
\label{thm:polar-coordinates}
\dcref{appC:thm-4}
\group{Coarea Formula}
(i) Let $f : \mathbb{R}^n \to \mathbb{R}$ be continuous and summable. Then
$$\int_{\mathbb{R}^n} f \, dx = \int_0^\infty \left( \int_{\partial B(x_0, r)} f \, dS \right) dr$$
for each point $x_0 \in \mathbb{R}^n$.

(ii) In particular
$$\frac{d}{dr} \left( \int_{B(x_0, r)} f \, dx \right) = \int_{\partial B(x_0, r)} f \, dS$$
for each $r > 0$.
\end{theorem}

Theorem \ref{thm:polar-coordinates} is a special case of

\begin{theorem}[Coarea formula]
\label{thm:coarea-formula}
\dcref{appC:thm-5}
\group{Coarea Formula}
Let $u : \mathbb{R}^n \to \mathbb{R}$ be Lipschitz continuous and assume that for a.e.\ $r \in \mathbb{R}$ the level set
$$\{x \in \mathbb{R}^n \mid u(x) = r\}$$
is a smooth, $(n-1)$-dimensional hypersurface in $\mathbb{R}^n$. Suppose also $f : \mathbb{R}^n \to \mathbb{R}$ is continuous and summable. Then
$$\int_{\mathbb{R}^n} f|Du| \, dx = \int_{-\infty}^{\infty} \left(\int_{\{u=r\}} f \, dS\right) dr.$$
\end{theorem}

Theorem \ref{thm:polar-coordinates} follows from Theorem \ref{thm:coarea-formula} by taking $u(x) = |x - x_0|$. See \cite[Chapter 3]{EvansGariepy1992} for more on the coarea formula. The word ``coarea'' is pronounced, and sometimes spelled ``co-area''.

\subsection{Convolution and Smoothing}

We next introduce tools that will allow us to build smooth approximations to given functions.

\begin{definition}[Notation: $U_\epsilon$]
\label{def:epsilon-interior-set-notation}
\dcref{appC:C.4-def-1}
\group{Mollification}
If $U \subset \mathbb{R}^n$ is open, $\epsilon > 0$, write $U_\epsilon := \{x \in U \mid \dist(x, \partial U) > \epsilon\}$.
\end{definition}

\begin{definition}[Standard mollifier]
\label{def:standard-mollifier}
\dcref{appC:C.4-def-2}
\group{Mollification}
(i) Define $\eta \in C^\infty(\mathbb{R}^n)$ by
$$\eta(x) := \begin{cases} C \exp\left(\frac{-1}{|x|^2-1}\right) & \text{if } |x| < 1 \\ 0 & \text{if } |x| \geq 1, \end{cases}$$
the constant $C > 0$ selected so that $\int_{\mathbb{R}^n} \eta \, dx = 1$.

(ii) For each $\epsilon > 0$, set
$$\eta_\epsilon(x) := \frac{1}{\epsilon^n} \eta\left(\frac{x}{\epsilon}\right).$$

We call $\eta$ the \textit{standard mollifier}. The functions $\eta_\epsilon$ are $C^\infty$ and satisfy
$$\int_{\mathbb{R}^n} \eta_\epsilon \, dx = 1, \quad \spt(\eta_\epsilon) \subset B(0,\epsilon).$$
\end{definition}

\begin{definition}[Mollification]
\label{def:mollification}
\dcref{appC:C.4-def-3}
\group{Mollification}
\uses{def:standard-mollifier,def:epsilon-interior-set-notation}
If $f : U \to \mathbb{R}$ is locally integrable, define its mollification
$$f^\epsilon := \eta_\epsilon * f \quad \text{in } U_\epsilon.$$

That is,
$$f^\epsilon(x) = \int_U \eta_\epsilon(x-y)f(y) \, dy = \int_{B(0,\epsilon)} \eta_\epsilon(y)f(x-y) \, dy$$
for $x \in U_\epsilon$.
\end{definition}

\begin{theorem}[Properties of mollifiers]
\label{thm:properties-of-mollifiers}
\dcref{appC:thm-6}
\group{Mollification}
\uses{def:mollification,def:standard-mollifier}
(i) $f^\epsilon \in C^\infty(U_\epsilon)$.

(ii) $f^\epsilon \to f$ a.e.\ as $\epsilon \to 0$.

(iii) If $f \in C(U)$, then $f^\epsilon \to f$ uniformly on compact subsets of $U$.

(iv) If $1 \le p < \infty$ and $f \in L^p_{\text{loc}}(U)$, then $f^\epsilon \to f$ in $L^p_{\text{loc}}(U)$.
\end{theorem}

\begin{proof}
1. Fix $x \in U_\epsilon$, $i \in \{1, \ldots, n\}$, and $h$ so small that $x + h e_i \in U_\epsilon$. Then

$$\frac{f^\epsilon(x + h e_i) - f^\epsilon(x)}{h} = \frac{1}{\epsilon^n} \int_U \frac{1}{h} \left[\eta\left(\frac{x + h e_i - y}{\epsilon}\right) - \eta\left(\frac{x - y}{\epsilon}\right)\right] f(y) \, dy$$

$$= \frac{1}{\epsilon^n} \int_V \frac{1}{h} \left[\eta\left(\frac{x + h e_i - y}{\epsilon}\right) - \eta\left(\frac{x - y}{\epsilon}\right)\right] f(y) \, dy$$

for some open set $V \subset\subset U$. As

$$\frac{1}{h}\left[\eta\left(\frac{x + h e_i - y}{\epsilon}\right) - \eta\left(\frac{x - y}{\epsilon}\right)\right] \to \frac{1}{\epsilon} \frac{\partial \eta}{\partial x_i}\left(\frac{x - y}{\epsilon}\right)$$

uniformly on $V$, $\frac{\partial f^\epsilon}{\partial x_i}(x)$ exists and equals

$$\int_U \frac{\partial \eta_\epsilon}{\partial x_i}(x - y) f(y) \, dy.$$

A similar argument shows that $D^\alpha f^\epsilon(x)$ exists, and

$$D^\alpha f^\epsilon(x) = \int_U D^\alpha \eta_\epsilon(x - y) f(y) \, dy \quad (x \in U_\epsilon),$$

for each multiindex $\alpha$. This proves (i).

2. According to Lebesgue's Differentiation Theorem (\S E.4),

$$(3) \quad \lim_{r \to 0} \int_{B(x,r)} |f(y) - f(x)| \, dy = 0$$

for a.e.\ $x \in U$. Fix such a point $x$. Then

$$|f^\epsilon(x) - f(x)| = \left|\int_{B(x,\epsilon)} \eta_\epsilon(x - y)[f(y) - f(x)] \, dy\right|$$

$$\le \frac{1}{\epsilon^n} \int_{B(x,\epsilon)} \eta\left(\frac{x - y}{\epsilon}\right) |f(y) - f(x)| \, dy$$

$$\le C \int_{B(x,\epsilon)} |f(y) - f(x)| \, dy \to 0 \quad \text{as } \epsilon \to 0,$$

by (3). Assertion (ii) follows.

3. Assume now $f \in C(U)$. Given $V \subset\subset U$ we choose $V \subset\subset W \subset\subset U$ and note that $f$ is uniformly continuous on $W$. Thus the limit (3) holds uniformly for $x \in V$. Consequently the calculation above implies $f^\epsilon \to f$ uniformly on $V$.

4. Next, assume $1 \leq p < \infty$ and $f \in L^p_{\text{loc}}(U)$. Choose an open set $V \subset\subset U$ and, as above, an open set $W$ so that $V \subset\subset W \subset\subset U$. We claim that for sufficiently small $\epsilon > 0$

$$(4) \qquad \|f^\epsilon\|_{L^p(V)} \leq \|f\|_{L^p(W)}.$$

To see this, we note that if $1 < p < \infty$ and $x \in V$,

$$|f^\epsilon(x)| = \left|\int_{B(x,\epsilon)} \eta_\epsilon(x-y)f(y)\,dy\right|$$

$$\leq \int_{B(x,\epsilon)} \eta_\epsilon^{1-1/p}(x-y)\eta_\epsilon^{1/p}(x-y)|f(y)|\,dy$$

$$\leq \left(\int_{B(x,\epsilon)} \eta_\epsilon(x-y)\,dy\right)^{1-1/p} \left(\int_{B(x,\epsilon)} \eta_\epsilon(x-y)|f(y)|^p\,dy\right)^{1/p}.$$

Since $\int_{B(x,\epsilon)} \eta_\epsilon(x-y)\,dy = 1$, this inequality implies

$$\int_V |f^\epsilon(x)|^p\,dx \leq \int_V \left(\int_{B(x,\epsilon)} \eta_\epsilon(x-y)|f(y)|^p\,dy\right) dx$$

$$\leq \int_W |f(y)|^p \left(\int_{B(y,\epsilon)} \eta_\epsilon(x-y)\,dx\right) dy = \int_W |f(y)|^p\,dy,$$

provided $\epsilon > 0$ is sufficiently small. This is inequality (4).

5. Now fix $V \subset\subset W \subset\subset U$, $\delta > 0$, and choose $g \in C(W)$ so that
$$\|f - g\|_{L^p(W)} < \delta.$$

Then
$$\|f^\epsilon - f\|_{L^p(V)} \leq \|f^\epsilon - g^\epsilon\|_{L^p(V)} + \|g^\epsilon - g\|_{L^p(V)} + \|g - f\|_{L^p(V)}$$
$$\leq 2\|f - g\|_{L^p(W)} + \|g^\epsilon - g\|_{L^p(V)} \text{ by (4)}$$
$$\leq 2\delta + \|g^\epsilon - g\|_{L^p(V)}.$$

Since $g^\epsilon \to g$ uniformly on $V$, we have $\limsup_{\epsilon \to 0} \|f^\epsilon - f\|_{L^p(V)} \leq 2\delta$.
\end{proof}

\subsection{Inverse Function Theorem}

Let $U \subset \mathbb{R}^n$ be an open set and suppose $\mathbf{f} : U \to \mathbb{R}^n$ is $C^1$, $\mathbf{f} = (f^1, \ldots, f^n)$. Assume $x_0 \in U$, $z_0 = \mathbf{f}(x_0)$.

\begin{definition}[Notation: gradient matrix $D\mathbf{f}$]
\label{def:gradient-matrix-df-square-notation}
\dcref{appC:C.5-def-1}
\group{Inverse and Implicit Function Theorems}
\uses{def:vector-valued-function-notation}
Remember from \S A.4 that we write
$$D\mathbf{f} = \begin{pmatrix} f^1_{x_1} & \cdots & f^1_{x_n} \\ \vdots & \ddots & \vdots \\ f^n_{x_1} & \cdots & f^n_{x_n} \end{pmatrix}_{n \times n} = \text{gradient matrix of } \mathbf{f}.$$
\end{definition}

\begin{definition}[Jacobian $J\mathbf{f}$]
\label{def:jacobian-jf}
\dcref{appC:C.5-def-2}
\group{Inverse and Implicit Function Theorems}
\uses{def:gradient-matrix-df-square-notation}
$$J\mathbf{f} = \text{Jacobian of } \mathbf{f} = |\det D\mathbf{f}| = \left| \frac{\partial(f^1, \ldots, f^n)}{\partial(x_1, \ldots, x_n)} \right|.$$
\end{definition}

\begin{figure}[h]\centering
% [FIGURE: Diagram showing an open set U (shaded oval region) in ℝⁿ containing a point x₀, with a function f mapping to a point z₀ in ℝⁿ]
\end{figure}

\begin{theorem}[Inverse Function Theorem]
\label{thm:inverse-function-theorem}
\dcref{appC:thm-7}
\group{Inverse and Implicit Function Theorems}
\uses{def:jacobian-jf}
Assume $\mathbf{f} \in C^1(U; \mathbb{R}^n)$ and
$$J\mathbf{f}(x_0) \neq 0.$$

Then there exist an open set $V \subset U$, with $x_0 \in V$, and an open set $W \subset \mathbb{R}^n$, with $z_0 \in W$, such that

(i) the mapping
$$\mathbf{f} : V \to W$$
is one-to-one and onto, and

(ii) the inverse function
$$\mathbf{f}^{-1} : W \to V$$
is $C^1$.

(iii) If $\mathbf{f} \in C^k$, then $\mathbf{f}^{-1} \in C^k$ ($k = 2, \ldots$).
\end{theorem}

\begin{figure}[h]\centering
% [FIGURE: Two regions side by side. Left region in $\mathbb{R}^n$ shows an open set with a subset V marked and point $x_0$ indicated. Right region in $\mathbb{R}^m$ shows a region with point $z_0$ marked. Arrows labeled $f$ point from left region to right region, and arrows labeled $f^{-1}$ point from right back to left. The domains are labeled $\mathbb{R}^n$ and $\mathbb{R}^m$ respectively.]
\end{figure}

\subsection{Implicit Function Theorem}

Let $n, m$ be positive integers.

\begin{definition}[Notation: points in $\mathbb{R}^{n+m}$]
\label{def:xy-point-notation-rnm}
\dcref{appC:C.6-def-1}
\group{Inverse and Implicit Function Theorems}
We write a typical point in $\mathbb{R}^{n+m}$ as
$$(x, y) = (x_1, \ldots, x_n, y_1, \ldots, y_m)$$
for $x \in \mathbb{R}^n$, $y \in \mathbb{R}^m$.
\end{definition}

Let $U \subset \mathbb{R}^{n+m}$ be an open set and suppose $\mathbf{f}: U \to \mathbb{R}^m$ is $C^1$, $\mathbf{f} = (f^1, \ldots, f^m)$. Assume $(x_0, y_0) \in U$, $z_0 = \mathbf{f}(x_0, y_0)$.

\begin{definition}[Notation: gradient matrix $D\mathbf{f}$]
\label{def:gradient-matrix-df-rectangular-notation}
\dcref{appC:C.6-def-2}
\group{Inverse and Implicit Function Theorems}
\uses{def:xy-point-notation-rnm}
$$D\mathbf{f} = \begin{pmatrix} f^1_{x_1} & \cdots & f^1_{x_n} & f^1_{y_1} & \cdots & f^1_{y_m} \\ \vdots & & \vdots & \vdots & & \vdots \\ f^m_{x_1} & \cdots & f^m_{x_n} & f^m_{y_1} & \cdots & f^m_{y_m} \end{pmatrix}_{m \times (n+m)}$$
$$= (D_x \mathbf{f}, D_y \mathbf{f}) = \text{gradient matrix of } \mathbf{f}.$$
\end{definition}

\begin{definition}[Jacobian $J_y\mathbf{f}$]
\label{def:jacobian-jyf}
\dcref{appC:C.6-def-3}
\group{Inverse and Implicit Function Theorems}
\uses{def:gradient-matrix-df-rectangular-notation}
$$J_y \mathbf{f} = |\det D_y \mathbf{f}| = \left| \frac{\partial(f^1, \ldots, f^m)}{\partial(y_1, \ldots, y_m)} \right|.$$
\end{definition}

\begin{figure}[h]\centering
% [FIGURE: Two regions side by side. Left region shows an open set $U$ in $\mathbb{R}^{n+m}$, drawn as an oval containing a marked point $(x_0,y_0)$ lying above a marked point $x_0$ on a coordinate axis labeled $\mathbb{R}^n$. Right region shows a line labeled $\mathbb{R}^m$ with a marked point $z_0$. A curved arrow labeled $f$ points from the oval in $\mathbb{R}^{n+m}$ to the point $z_0$ in $\mathbb{R}^m$.]
\end{figure}

\begin{theorem}[Implicit Function Theorem]
\label{thm:implicit-function-theorem}
\dcref{appC:thm-8}
\group{Inverse and Implicit Function Theorems}
\uses{def:jacobian-jyf}
Assume $\mathbf{f} \in C^1(U; \mathbb{R}^m)$ and
$$J_y \mathbf{f}(x_0, y_0) \neq 0.$$

Then there exists an open set $V \subset U$, with $(x_0, y_0) \in V$, an open set $W \subset \mathbb{R}^n$, with $x_0 \in W$, and a $C^1$ mapping $\mathbf{g} : W \to \mathbb{R}^m$ such that

(i) $\mathbf{g}(x_0) = y_0$,

(ii) $\mathbf{f}(x, \mathbf{g}(x)) = z_0 \quad (x \in W)$,

and

(iii) if $(x,y) \in V$ and $\mathbf{f}(x,y) = z_0$, then $y = \mathbf{g}(x)$.

(iv) If $\mathbf{f} \in C^k$, then $\mathbf{g} \in C^k$ ($k = 2, \ldots$).

The function $\mathbf{g}$ is implicitly defined near $x_0$ by the equation $\mathbf{f}(x,y) = z_0$.
\end{theorem}

\begin{figure}[h]\centering
% [FIGURE: The same two-region diagram as above (open set $U \subset \mathbb{R}^{n+m}$ containing shaded subset $V$ around $(x_0,y_0)$, mapping via $f$ to $z_0 \in \mathbb{R}^m$), now with an additional shaded region in $\mathbb{R}^n$ around a point $x_0$ and a downward-then-across arrow labeled $g$ running from that shaded region in $\mathbb{R}^n$ up into the shaded subset $V$, illustrating the implicitly defined function $g$.]
\end{figure}

\subsection{Uniform Convergence}

We record here the Arzel\`a-Ascoli compactness criterion for uniform convergence:

\begin{theorem}[Arzel\`a--Ascoli Theorem]
\label{thm:arzela-ascoli-theorem}
\dcref{appC:thm-9}
\group{Uniform Convergence}
\uses{def:uniform-equicontinuity}
Suppose that $\{f_k\}_{k=1}^{\infty}$ is a sequence of real-valued functions defined on $\mathbb{R}^n$, such that
$$|f_k(x)| \leq M \quad (k = 1, \ldots, x \in \mathbb{R}^n)$$
for some constant $M$, and the $\{f_k\}_{k=1}^{\infty}$ are \textit{uniformly equicontinuous}. Then there exists a subsequence $\{f_{k_j}\}_{j=1}^{\infty} \subseteq \{f_k\}_{k=1}^{\infty}$ and a continuous function $f$, such that
$$f_{k_j} \to f \quad \text{uniformly on compact subsets of } \mathbb{R}^n.$$
\end{theorem}

\begin{definition}[Uniform equicontinuity]
\label{def:uniform-equicontinuity}
\dcref{appC:def-9}
\group{Uniform Convergence}
To say the $\{f_k\}_{k=1}^{\infty}$ are uniformly equicontinuous means that for each $\varepsilon > 0$, there exists $\delta > 0$, such that $|x - y| < \delta$ implies $|f_k(x) - f_k(y)| < \varepsilon$, for $x, y \in \mathbb{R}^n$, $k = 1, \ldots$
\end{definition}
